\section{Discussion}
\label{sec:discussion}

\paragraph{Why failure-trace gating is effective.}
The two gates in FTA target failure modes that are identifiable at runtime
without any access to gold labels.
The Low-Depth Guard (LDG) addresses a structural weakness: when the frontier
search terminates with a single weak branch---signalled by the logged
override-reason field---the resulting answer is drawn from a narrow slice of
the candidate space.
In such cases, the external majority vote is empirically more reliable.
The External-Consensus Fallback (ECO) targets a complementary pattern: when
the frontier selected its answer with high internal confidence, yet all three
independent external baselines unanimously prefer a different answer, that
unanimous disagreement provides strong evidence that the frontier selection
is wrong.
Both signals are available as logged metadata from normal frontier execution;
no additional inference or labelling is required.

\paragraph{When FTA helps and when it does not.}
FTA provides the greatest benefit on examples where the frontier's failure mode
is structurally signalled---narrow search depth or unanimous external
disagreement---because these are precisely the cases the two gates are designed
to detect.
In the Final-300 evaluation, LDG fires on 63 of 300 examples (21.0\%) and
yields 36 wins against 9 losses versus the frontier, while ECO fires on only
3 examples (1.0\%) with zero losses.
The no-gate majority (234/300, 78.0\%) retains the frontier answer with 94.4\%
accuracy, indicating that the frontier is reliable on most examples and the
policy correctly defers to it.

FTA does not help---and cannot---when all methods in the candidate pool produce
incorrect answers for the same example.
This \emph{pool-miss} failure class accounts for 24 of the 40 residual errors
on Final-300 (60\%), representing an irreducible floor for any post-generation
selection policy operating over the current pool.
Additionally, the ECO gate requires unanimous external agreement, so cases
where the external baselines partially but not unanimously disagree with the
frontier are handled by the frontier default.
Extending the method to handle partial agreement would require stronger
statistical justification, as such cases showed more ambiguous outcomes in
development.

The current evaluation is scoped to GSM8K under the Cohere
command-r-plus-08-2024 setting.
Preliminary evidence on other providers (Section~\ref{sec:limitations}) suggests
the failure-trace signal may have broader utility, but the main findings should
not be extrapolated to other benchmarks or providers without independent
replication.

\paragraph{Not-retained variants.}
Several algorithmic extensions were developed and evaluated during the research
but did not meet validation-grade evidence criteria
(Table~\ref{tab:negative_results}).
A TALE-default router produced zero net switching effect across evaluation
seeds, equalling TALE performance exactly on seeds~41 and~61, because the
assumed agreement region did not materialise consistently.
An extra-action variant---collecting additional frontier or retry actions for
predicted failure cases---showed a non-negative initial signal in a 40-case
pilot but failed to replicate on a disjoint 80-case follow-up
(net $-1$ for extra frontier actions; net $-4$ for TALE retries), suggesting
the pilot signal was run-specific rather than generalisable.
A cluster-level selector yielded only limited offline gain without sufficient
aggregate validation.
These outcomes indicate that rule-based post-generation selection gains within
the current pool composition are likely close to saturation, and further
improvements will require better candidate generation or learned allocation
policies.

\paragraph{Evidence from the Aggregate-720 evaluation.}
The most substantive evidence shift in this work is the transition from a
single-run 300-example evaluation (seed~41, where the CI lower bound vs.\ L1
crossed zero at $[-0.48,\, +5.71]$) to the disjoint Aggregate-720 evaluation
(seeds 41, 61, 71).
By combining three independently sampled evaluation splits, the
source-stratified bootstrap CI narrows sufficiently that the lower bound
versus every external baseline is strictly positive.
The source-stratified design controls for seed-specific variance, and the
zero-overlap check on example-id and question-hash confirms independence.
The point-estimate advantage of FTA is consistent across all three seeds,
which is the primary reason the aggregate CIs no longer include zero.

\paragraph{Multi-provider supporting evidence.}
As supporting evidence for the failure-trace signal,
a preliminary multi-provider experiment (D9) trains a learned gate on the same
failure-trace features across three providers.
Cross-validated accuracy on the D9 gate reaches $0.5018 \pm 0.0252$ against a
frontier baseline of $0.3436$, with zero false-override events across all tested
providers and decision thresholds.
The Cloudrift/Qwen provider contributes the strongest rescue signal
(frontier 37.5\%, D6-gated 55.0\%, 24 unique-correct outcomes, zero
regressions with lenient extraction).
These results are suggestive rather than conclusive about broader utility, but
the evaluation scope is preliminary (80--160 examples per provider); the
findings constitute supporting evidence, not a second primary claim.
Broader multi-provider evaluation at matched scale is a natural next step.

\paragraph{Evidence scope and generalization.}
The main findings are scoped to the tested evaluation setting.
What this evidence does \emph{not} establish is generalization beyond that
scope: whether the same failure-trace signals are informative on other
providers, benchmarks, or budget levels is an open empirical question.
The most efficient next replication would be a second GSM8K provider run or a
same-provider second benchmark with the identical frontier, L1, S1, and TALE
logging contract.
Section~\ref{sec:limitations} enumerates the scope boundaries explicitly.
